\originalsection{场模态如何成为谐振子}
先把分析力学的关系说明: 对坐标 $q_\alpha$, 拉格朗日量 $L(q,\dot q)$ 给 $p_\alpha=\partial L/\partial\dot q_\alpha$, 哈密顿量 $H=\sum p_\alpha\dot q_\alpha-L$. 自由电磁场在真空的拉格朗日量为
$L=\int[\eps E^2/2-B^2/(2\mu_0)]\dd^3x$.
这与能量 $U=\int[\eps E^2/2+B^2/(2\mu_0)]\dd^3x$ 的磁项符号不同; 前者产生运动方程, 后者衡量能量.

在无静电背景的无源腔取 $V=0$, 用 $\vct A=\sum q_\alpha\vct f_\alpha$, 就有 $\ee=-\sum\dot q_\alpha\vct f_\alpha$、$\bb=\sum q_\alpha\nabla\times\vct f_\alpha$. 导体边界选择 $\nn\times\vct f_\alpha=0$, 模态横向且满足 $\nabla\times\nabla\times\vct f_\beta=\omega_\beta^2\vct f_\beta/c^2$. 叉积散度恒等式及边界条件给
\[
 \int(\nabla\times\vct f_\alpha)\cdot(\nabla\times\vct f_\beta)\dd^3x
 =\frac{\omega_\beta^2}{c^2}\int\vct f_\alpha\cdot\vct f_\beta\dd^3x.
\]
于是 $\eps\int\vct f_\alpha\cdot\vct f_\beta=\delta_{\alpha\beta}$ 消去交叉项, 得 $L=\sum(\dot q_\alpha^2-\omega_\alpha^2q_\alpha^2)/2$.
欧拉--拉格朗日方程为 $\ddot q_\alpha+\omega_\alpha^2q_\alpha=0$, $p_\alpha=\dot q_\alpha$, $H=\sum(p_\alpha^2+\omega_\alpha^2q_\alpha^2)/2$. 模态不是实际带电小球, 其“坐标”是场的展开振幅. 这一推导给研究生作业中的腔模正交及场量子化入口, 没有把规范纵向自由度当成额外物理振子.

\originalsection{带源场作用量与因果模态解}
8.311补充笔记\textit{Field Oscillators}第1--2页将自由场模态推广到带源场. 以下用SI制重新推导, 沿用上一节$\eps\int\vct f_\alpha\cdot\vct f_\beta\dd^3x=\delta_{\alpha\beta}$的归一, 不照搬原笔记高斯制中的$c$和$4\pi$系数. 先在全空间真空中取平滑局域源$\rho,\jj$, 满足$\partial_t\rho+\nabla\cdot\jj=0$, 场的远处衰减足以消去分部积分边界项. 带源拉格朗日量及作用量为
\[
L=\int\left[\frac{\eps}{2}E^2-\frac{B^2}{2\mu_0}+\jj\cdot\vct A-\rho V\right]\dd^3x,
\qquad S=\int L\dd t.
\]
源在这里给定; 若同时求粒子轨道, 还须加入粒子的机械作用量. $\jj\cdot\vct A-\rho V$是源与势的耦合, 不能把自由场能量直接当成这一带源作用量.

取库仑规范$\nabla\cdot\vct A=0$, $\ee=-\dot{\vct A}-\nabla V$. 电场平方的交叉项消失, 因为
$\int\dot{\vct A}\cdot\nabla V\dd^3x=-\int V\nabla\cdot\dot{\vct A}\dd^3x=0$.
关于$V$的变分因而为
$\delta_VL=\int(-\eps\nabla^2V-\rho)\delta V\dd^3x$,
给$\nabla^2V=-\rho/\eps$. $V$没有时间导数项, 是每个时刻由电荷约束确定的量, 不是额外的传播振子. 取$V\to0$的边界条件, 则
$V(\rr,t)=(4\pi\eps)^{-1}\int\rho(\rr',t)/|\rr-\rr'|\dd^3x'$.

把该约束解代回作用量, 必须同时保留纵向场能与$-\rho V$耦合. 分部积分给
\[
\frac{\eps}{2}\int|\nabla V|^2\dd^3x
=-\frac{\eps}{2}\int V\nabla^2V\dd^3x
=\frac12\int\rho V\dd^3x,
\]
所以约化结果为
\begin{align*}
L_{\rm red}&=\int\left[\frac{\eps}{2}|\dot{\vct A}|^2
-\frac{|\nabla\times\vct A|^2}{2\mu_0}+\jj\cdot\vct A\right]\dd^3x-U_C(t),\\
U_C(t)&=\frac12\int\rho V\dd^3x
=\frac1{8\pi\eps}\iint\frac{\rho(\rr,t)\rho(\rr',t)}{|\rr-\rr'|}\dd^3x\dd^3x'.
\end{align*}
库仑项在拉格朗日量中是$-U_C$, 二分之一来自代回约束后的两项合并. 本式先对有限能量的平滑源成立; 点电荷的自身项发散, 不能把它直接算作有限相互作用能. 不同点荷的相互作用部分才是$\sum_{i<j}q_iq_j/(4\pi\eps r_{ij})$. 有界导体腔须将库仑核改成满足壁面条件的空间Green核. 此处同一时刻的$V$体现库仑规范的约束性质; 可观测电场还含$-\dot{\vct A}$, 不能仅据$V$判断信号传播.

继续用横向模态展开$\vct A=\sum_\alpha q_\alpha\vct f_\alpha$. 取上一节的正交归一和边界条件（全空间可先作大盒离散化再取极限）, 定义$F_\alpha(t)=\int\jj(\rr,t)\cdot\vct f_\alpha(\rr)\dd^3x$. 模态只耦合电流的横向部分: 任意可消去边界项的梯度电流$\nabla\psi$与$\vct f_\alpha$的内积为零. 因而
$L_{\rm red}=\sum_\alpha[\dot q_\alpha^2/2-\omega_\alpha^2q_\alpha^2/2+q_\alpha F_\alpha]-U_C$,
对$q_\alpha$变分给$\ddot q_\alpha+\omega_\alpha^2q_\alpha=F_\alpha(t)$, 没有额外的$\eps$或$c$因子.

对$\omega_\alpha>0$及连续驱动力, 在$t_0$取零初值$q_\alpha(t_0)=\dot q_\alpha(t_0)=0$. 因果Green函数为$G_\alpha(t-s)=\Theta(t-s)\sin[\omega_\alpha(t-s)]/\omega_\alpha$, 解为
\[
q_\alpha(t)=\int_{t_0}^{t}\frac{\sin[\omega_\alpha(t-s)]}{\omega_\alpha}F_\alpha(s)\dd s,
\qquad t\ge t_0.
\]
它只用过去的驱动力. 在$t=t_0$积分为零; 对$t$求导时上端的正弦为零, 得$\dot q_\alpha=\int_{t_0}^t\cos[\omega_\alpha(t-s)]F_\alpha(s)\dd s$, 初始导数也为零. 再求导, 上端余弦为1, 故$\ddot q_\alpha=F_\alpha(t)-\omega_\alpha^2q_\alpha$, 验证了运动方程. 初值非零时另加$q_\alpha(t_0)\cos[\omega_\alpha(t-t_0)]+\dot q_\alpha(t_0)\sin[\omega_\alpha(t-t_0)]/\omega_\alpha$. 若边界允许零频模, 核的连续极限是$t-s$; 从$-\infty$写零入射解还需保证积分收敛, 或给出明确绝热极限.

作为点源耦合的分布极限, $\jj=q\vct v(t)\delta^3(\rr-\rr_p(t))$给$F_\alpha=q\vct v\cdot\vct f_\alpha(\rr_p)$; 它不消除点源库仑自能的限制. 将正弦、余弦模的实际SI归一代入此式, 就得到后文运动电荷驱动场模及Frank--Tamm共振计算的入口. 因果卷积适用于一般时间驱动, 后文的余弦特解只是它的一个可显式积分的情形.

\originalsection{由四势变换求电磁场变换}
相对论先修只需这里的洛伦兹坐标变换和四矢量规则. 对沿 $+x$ 运动的新系, 四势 $(V/c,\vct A)$ 按四矢量变换:
$V'=\gamma(V-vA_x)$、$A_x'=\gamma(A_x-vV/c^2)$, $A_y'=A_y$, $A_z'=A_z$.
逆坐标关系的链式法则给 $\partial_{t'}=\gamma(\partial_t+v\partial_x)$、$\partial_{x'}=\gamma(\partial_x+v\partial_t/c^2)$. 例如
\begin{align*}
 E_y'&=-\partial_yV'-\partial_{t'}A_y'
 =\gamma[E_y+v(\partial_yA_x-\partial_xA_y)]
 =\gamma(E_y-vB_z),\\
 B_z'&=\partial_{x'}A_y'-\partial_yA_x'
 =\gamma[B_z+(v/c^2)(\partial_tA_y+\partial_yV)]
 =\gamma(B_z-vE_y/c^2).
\end{align*}
同样得到 $E_z'=\gamma(E_z+vB_y)$、$B_y'=\gamma(B_y+vE_z/c^2)$; 平行分量的两次 $\gamma$ 乘 $1-v^2/c^2$ 给 $E_x'=E_x$, $B_x'=B_x$. 这说明电磁场如何因更换观察系而混合, 并解释式中速度的符号由“新系相对旧系”决定.

不变量也可代入验证, 无需只背结论. 例如
$(E_y')^2-c^2(B_z')^2=\gamma^2(1-v^2/c^2)(E_y^2-c^2B_z^2)$,
另一横向对相同, 平行分量不变, 所以 $E'^2-c^2B'^2=E^2-c^2B^2$.
对点积, 两个横向对的混合项互相抵消, 得
$E_y'B_y'+E_z'B_z'=\gamma^2(1-v^2/c^2)(E_yB_y+E_zB_z)$,
加上 $E_xB_x$ 即得 $\ee'\cdot\bb'=\ee\cdot\bb$. 这给研究生作业9寻找纯电或纯磁参考系的工具, 但不意味着所有电磁场都能这样消去一种场; 特别是两不变量都为零的非零平面波没有静止系.

\originalsection{介质中运动电荷的模态共振与Frank--Tamm谱}
8.311的四页\texttt{charge\_mov.pdf}用受迫场振子解释切伦科夫与碰撞辐射. 先取透明、无色散、非磁介质$\varepsilon=\eps n^2$, $n>1$, 在体积$\mathcal V$的大盒中将横势按实正弦/余弦模展开. 模函数归一为$\vct f_{\vct k,s}=\sqrt{2/(\varepsilon\mathcal V)}\vct e_s\cos(\vct k\cdot\rr)$或正弦, $\vct e_s\perp\vct k$, 固有频率$\omega_k=ck/n$. 粒子耦合$q\vct v\cdot\vct A$使每模
$\ddot Q_{ks}+\omega_k^2Q_{ks}=F_{ks}(t)$, $F_{ks}=q\vct v\cdot\vct f_{ks}(\rr_q(t))$.

匀速$\rr_q=\vct vt$驱动力频率$\omega_f=kv\cos\theta$. 以余弦驱动$F_a\cos\omega_ft$且初始场模为零, 实解$Q=F_a(\cos\omega_ft-\cos\omega_kt)/(\omega_k^2-\omega_f^2)$. 共振附近累计能的增长率在长时间极限给$\pi F_a^2\delta(\omega_k-\omega_f)/4$, 由$\sin(\Delta t)/\Delta\to\pi\delta(\Delta)$得到. 正弦模贡献相同, 两横向偏振满足$\sum_s|\vct v\cdot\vct e_s|^2=v^2\sin^2\theta$. 实模只对独立半球波矢计数, 相反$\vct k$不是另一个实模式.

以$\cos\theta\ge0$选半球, 模态密度$\mathcal V\dd^3k/(2\pi)^3$给
\[
P=\frac{q^2}{8\pi^2\varepsilon}\int_{\rm half}k^2\dd k\dd\Omega\,v^2\sin^2\theta\,\delta(ck/n-kv\cos\theta)
=\frac{q^2v}{4\pi\varepsilon}\int k\dd k\left(1-\frac{c^2}{n^2v^2}\right).
\]
只有$nv>c$时delta根在角区内. 换$k=n\omega/c$得正频谱$\dd P/\dd\omega=q^2v\omega[1-c^2/(n^2v^2)]/(4\pi\eps c^2)$, 或每路径长$\dd^2W/(\dd x\dd\omega)=q^2\omega[1-c^2/(n^2v^2)]/(4\pi\eps c^2)$. 这是非磁透明介质的Frank--Tamm式, 不是匀速电荷在真空辐射; 在真空亚光速根根本不存在. 真实材料应按$n(\omega)$频带积分并检查吸收及色散储能, 不把常$n$谱无限外推.

碰撞时切换速度的驱动模有两种频率, 连续匹配$Q,\dot Q$后留下自由振子辐射; 第10章的角谱推导已将它等价写成两端速度软因子. 分母$1-\nn\cdot\vct v/c$使高速度辐射集中于运动方向. 软谱趋常量所以低频辐射能可积, 但若另用量子能$\hbar\omega$计光子数, $\int(\dd W/\dd\omega)\dd\omega/(\hbar\omega)$形式上对数发散. 原笔记据此提到红外问题; 这里只解释经典能谱与光子计数的差别, 不冒称完成量子电动力学红外消除证明.
